%!TEX root = ../../main.tex
\chapter{B 需求}
% ============================================================

\section{系统用例}

用例（Use Case）即使用实例，从用户的角度出发描述计算机能够实现的具体目标。
系统用例图从改进业务序列图映射而来：系统用例 = 从外部指向目标系统的消息，
系统执行者 = 与目标系统实例有实线相连的实例所代表的系统，没有辅执行者。
其中登记成交、登记退货、登记货品入库、预定商品、预约查看货品、查看调货价格、
提交货品出让申请、审核出让申请、预约上门时间、查看加工任务、提交完工交付申请、
登记货款记录与维护调货记录 13 个用例直接来自六个改进业务场景的消息映射（见 A 部分改进业务序列图一节）；
登记寄售合同、结算寄售货款、结算调货台账、登记分期计划、记录分期还款与查看库存流水 6 个用例
来自新增的供应商寄售场景（图~\ref{fig:15-to-be-consignment.png}）与台账、分期、流水三项扩展；
维护客户信息、登记送检与证书、查看销售与周转统计、维护基础数据、登录、浏览货品、
开通外部账户 7 个用例是支撑上述用例与愿景指标所必需的补充；
查询预约与跟进预定到店 2 个用例为预约环节补上凭证找回与到店跟进。
系统的执行者分为老板、销售、客户、合作商户、供应商与加工师傅六类，
故按执行者端进行用例分析。总体用例图见图 \ref{fig:08-system-usecase.png}，
共 6 个执行者、28 个用例；
系统负责货品档案、成交记录、统计，面向客户的货品展示、预约与预定，
以及供应商与加工师傅的在线协同，不接入在线支付接口，不参与鉴定和物流。

\figab[0.60]{系统用例图：系统边界、执行者与用例}{08-system-usecase.png}

账号角色设老板、销售、合作商户、供应商与加工师傅五种，
合作商户、供应商与加工师傅的账号都由老板逐个开通（UC-12），不开放自助注册；
客户浏览、预约与预定无需账号，前端不展示价格——标出的价格并不是最低售价。
柜台端用电脑操作；供应商、加工师傅与同行商户多在散档口、不便用电脑且文化程度不一，
外部协同端做成微信小程序，操作步骤尽量少，账号与初始口令由老板线下交付。
入库与出库在现实里由老板和销售同时经手，改进后的口径是默认分工加系统兜底：
出库（成交、调出、退货）由销售登记，入库建档以老板为主，
老板不在店时由销售代理收货建档，系统保留经手人，记错能追到责任人；
送检由销售在系统提交货品清单并每天去鉴定机构集中办理，同样留经手人。
线上预约与上门预约统一排期：客户在系统约好到店时段后直接找对应销售，
避免销售当天已有接待任务时老板两头协调。
“登录”本身不产生业务价值，但除顾客端外，进入各端功能都要先验证身份，
单独列为用例（UC-08），由各端主用例以 \code{<<include>>} 指向。

按执行者端划分的用例如下。

\subsection{管理端用例（老板）}

管理端由老板使用。老板在系统中承担经营决策者的角色，拥有登录、登记货品入库、
查看销售与周转统计、维护基础数据、开通外部账户、审核出让申请、登记货款记录、
维护调货记录、登记寄售合同、结算寄售货款、结算调货台账与查看库存流水等用例，
具体用例图如图所示。

\figab[0.88]{管理端用例}{uc-boss.png}

用例详细说明如下：

\noindent（1）登录：老板使用由自己开通或设置的账号登录，验证通过后进入管理端功能。

\noindent（2）登记货品入库：对采购、回收、调货调入、加工完工与寄售的货品逐件建档，
录入来源、成本与最低售价，货品状态置为“在柜”，周转周期从这里开始计时。

\noindent（3）查看销售与周转统计：系统按成交记录自动累计月度销售额、销售客户数
与货品加权周转周期，老板随时查看，是愿景指标的读数入口。

\noindent（4）维护基础数据：维护来源方（供应商/加工师傅/同行商户）、鉴定机构与货品类别，
支撑其余用例的参照数据。

\noindent（5）开通外部账户：为合作商户、供应商与加工师傅逐个开通账号与权限，
停用即收回，外部角色不能自行注册。

\noindent（6）审核出让申请：审核供应商在线提交的货品出让申请，
通过的转入入库流程，驳回的填写原因并留痕。

\noindent（7）登记货款记录：登记采购货款、调货货款、加工工费与寄售货款的新增与收付，
每笔货款可追溯到来源单据与经手人。

\noindent（8）维护调货记录：维护调货往来记录与调货库存记录，
调货货品去向可追溯，结算不扯皮。

\noindent（9）登记寄售合同：供应商寄存货品凭合同建档入库，
寄售价记为成本作为结算依据，合同与明细一次登记。

\noindent（10）结算寄售货款：已售出的寄售货品按寄售价与供应商结算，
明细置为已结算，回款有结算单为据。

\noindent（11）结算调货台账：按商户汇总未结算调货记录一次对账，
结算后单据置为已结算，鉴赏期届满即结清。

\noindent（12）查看库存流水：任一货品的每次状态变更都留有流水，
进出有账可查，是调货库存一本账的追溯底账。

\subsection{柜台端用例（销售）}

柜台端由销售使用。销售在系统中承担柜台经手的角色，拥有登录、登记成交、
登记退货、维护客户信息、登记送检与证书、登记分期计划与记录分期还款等用例，
具体用例图如图所示。

\figab[0.88]{柜台端用例}{uc-sale.png}

用例详细说明如下：

\noindent（1）登录：销售使用老板开通的账号登录，验证通过后进入柜台端功能。

\noindent（2）登记成交：成交（全款或定金）当场登记到货品编号上，
系统校验货品在柜并更新货品状态，销售额与销售客户数两个度量都从这里来。

\noindent（3）登记退货：收回退货，货品状态改回“在柜”，冲减销售统计，
保证销售与库存数据准确。

\noindent（4）维护客户信息：维护客户姓名、联系方式、偏好与购买记录，
是精准营销的前提。

\noindent（5）登记送检与证书：登记送检记录与鉴定证书信息并挂到对应货品，
客户索证时随时可调，是高价货品销售的依据。

\noindent（6）登记分期计划：大额成交当场谈定分期，系统按成交价生成分期计划，
总额与期数一次确定。

\noindent（7）记录分期还款：录入每期还款并推进到期日，
结清自动归档，逾期计划随时可查。

\subsection{顾客端用例（客户）}

顾客端面向到店与线上客户。客户无需登录即可使用，拥有浏览货品、预约查看货品、
预定商品等用例，具体用例图如图所示。

\figab[0.88]{顾客端用例}{uc-cust.png}

用例详细说明如下：

\noindent（1）浏览货品：无需登录浏览在柜货品的图文展示，不显示价格，
把在柜货品摆到前端引流。

\noindent（2）预约查看货品：提交联系方式与到店时段，预约到店看货，转化到店。

\noindent（3）预定商品：留手机号预定在柜货品，一个手机号 3 次额度、
如约达成后恢复，锁定到店意向。

\subsection{商户端用例（合作商户）}

商户端由合作商户使用。合作商户拥有登录、查看调货价格等用例，具体用例图如图所示。

\figab[0.88]{商户端用例}{uc-merch.png}

用例详细说明如下：

\noindent（1）登录：合作商户使用老板开通的账号登录，验证通过后仅可见调货价目。

\noindent（2）查看调货价格：登录后查看在柜货品的调货价目，
减少问价与开户的人工往返。

\subsection{供应端用例（供应商）}

供应端由供应商使用（含玉料供应商、成品供应商与回流出让方）。供应商拥有登录、
提交货品出让申请、预约上门时间等用例，具体用例图如图所示。

\figab[0.88]{供应端用例}{uc-supp.png}

用例详细说明如下：

\noindent（1）登录：供应商使用老板开通的账号登录，验证通过后进入供应端功能。

\noindent（2）提交货品出让申请：在线提交玉料、成品或回流货品的出让申请，
货源线索集中在线上，不用上门先谈一轮。

\noindent（3）预约上门时间：申请通过后预约上门交接时段，按约前来，不白跑。

\subsection{加工端用例（加工师傅）}

加工端由加工师傅使用。加工师傅拥有登录、查看加工任务、提交完工交付申请、预约上门时间等用例，
具体用例图如图所示。

\figab[0.88]{加工端用例}{uc-work.png}

用例详细说明如下：

\noindent（1）登录：加工师傅使用老板开通的账号登录，验证通过后进入加工端功能。

\noindent（2）查看加工任务：登录查看分派给自己的加工任务、加工要求与工费，
不用电话逐个通知。

\noindent（3）提交完工交付申请：完工后在线报备，等待老板确认并安排交接入库，
完工信息第一时间可见。

\noindent（4）预约上门时间：在线预约上门交接时段，出让交货与完工交付都不用电话约定。

\section{用例清单}

\needspace{8\baselineskip}
\begin{table}[H]
\centering
\small
\caption{系统用例清单（排序依据是愿景）}
\begin{tabularx}{\textwidth}{@{}l l Y c@{}}
\toprule
编号 & 用例 & 说明 & 排序 \\
\midrule
UC-01 & 登记成交 & 成交（全款或定金）当场登记到货品编号上，系统校验并保存 & 1（关键） \\
UC-02 & 登记货品入库 & 采购、回收、调货、加工完工、寄售的货品逐件建档入库 & 2（关键） \\
UC-03 & 登记退货 & 收回退货，货品状态改回“在柜”，冲减销售统计 & 4 \\
UC-04 & 维护客户信息 & 客户姓名、联系方式、偏好与购买记录，支撑精准营销 & 5 \\
UC-05 & 登记送检与证书 & 送检记录与鉴定证书信息挂到对应货品 & 6 \\
UC-06 & 查看销售与周转统计 & 月度销售额、销售客户数、货品加权周转周期 & 3 \\
UC-07 & 维护基础数据 & 来源方（供应商/加工师傅/同行商户）、鉴定机构与货品类别 & 7 \\
UC-08 & 登录 & 验证用户名与口令后按角色进入系统（被 UC-01～UC-07 与 UC-11、UC-12、UC-14～UC-26、UC-28 包含） & — \\
UC-09 & 浏览货品 & 客户无需登录浏览在柜货品图文展示，不显示价格 & 8 \\
UC-10 & 预约查看货品 & 客户提交联系方式与到店时段，预约到店看货 & 9 \\
UC-11 & 查看调货价格 & 合作商户登录后查看在柜货品的调货价目 & 10 \\
UC-12 & 开通外部账户 & 老板为合作商户、供应商与加工师傅逐个开通账号与权限 & 11 \\
UC-13 & 预定商品 & 客户留手机号预定在柜货品，一个手机号 3 次额度、如约达成恢复 & 12 \\
UC-14 & 提交货品出让申请 & 供应商（含回流出让方）在线提交货品出让申请 & 13 \\
UC-15 & 预约上门时间 & 供应商与加工师傅在线预约上门交接时段 & 14 \\
UC-16 & 查看加工任务 & 加工师傅登录查看分派给自己的加工任务（含要求、工费与期限） & 15 \\
UC-17 & 提交完工交付申请 & 加工师傅完工后在线提交交付申请，等待老板确认并安排入库 & 16 \\
UC-18 & 审核出让申请 & 老板审核出让申请，通过后转入入库流程 & 17 \\
UC-19 & 登记货款记录 & 供货方货款、调货货款与加工工费的新增与收付记录 & 18 \\
UC-20 & 维护调货记录 & 调货往来记录与调货库存记录的新增与更新 & 19 \\
UC-21 & 登记寄售合同 & 供应商寄存货品逐件建档入库，寄售价作结算依据 & 20 \\
UC-22 & 结算寄售货款 & 已售寄售货品按寄售价与供应商结算，明细置已结算 & 21 \\
UC-23 & 结算调货台账 & 按商户汇总未结算调货记录，一次对账结清 & 22 \\
UC-24 & 登记分期计划 & 大额成交分期付款，总额等于成交价 & 23 \\
UC-25 & 记录分期还款 & 录入还款、推进到期日，结清自动归档 & 24 \\
UC-26 & 查看库存流水 & 货品每次状态变更留痕，进出有账可查 & 25 \\
UC-27 & 查询预约 & 客户凭手机号后四位查自己的预约，只显示编号、时段与货品 & 26 \\
UC-28 & 跟进预定到店 & 销售查看当日到店名单，电话联系后标记已联系并登记到店结论 & 27 \\
\bottomrule
\end{tabularx}
\end{table}

\noindent
逐项理由如下：UC-01 登记成交排第 1，销售额、销售客户数两个度量都从成交记录来；
UC-02 登记货品入库排第 2，周转周期从入库开始计时，没有入库记录的货品无法追踪；
UC-06 查看销售与周转统计排第 3，愿景指标的读数入口；
UC-03 退货排第 4，直接影响销售与库存数据的准确性；
UC-04 维护客户信息排第 5，精准营销的前提；UC-05 送检与证书排第 6，高价货品销售的依据；
UC-07 基础数据排第 7，支撑前述用例的参照数据；UC-08 被其他用例包含，不单独排序。
UC-09、UC-10 把在柜货品摆到前端引流并转化到店，排在台账功能之后；
UC-11、UC-12 服务同行调货，减少问价与开户的人工往返；
UC-13 预定商品锁定到店意向，UC-14～UC-18 把供应商、回流出让方与加工师傅的往来搬到线上，
省去电话问货与跑腿上门的时间；
UC-19、UC-20 承接改进场景里的记账消息，货款与调货一本账；
UC-21～UC-23 把寄售与台账结算搬上线，寄售价作结算依据、按商户一次对账，
钱账两清有据可查；UC-24～UC-25 承接大额成交的分期付款，逾期可追踪；
UC-26 库存流水给一本账补上追溯底账，货品每次状态变更留痕。

\section{系统用例规约}

根据各用例设计用例规约，进一步详细规定用例的用例编号与用例名、执行者、前置条件、
后置条件、涉众利益、基本路径、扩展，以及字段列表、业务规则、质量需求、设计约束
四个补充约束，每个规约用一张两列表格，下文按 UC-01～UC-26 逐个给出。
前置与后置条件只写系统能检测到的约束，扩展路径只写系统能感知并要处理的意外与分支。

\noindent（1）登记成交

{\small
\begin{longtable}{@{}L{2.4cm} L{12.9cm}@{}}
\caption{登记成交用例规约}\label{tab:uc-01}\\
\toprule
\textbf{标题} & \textbf{内容} \\
\midrule
\endfirsthead
\toprule
\textbf{标题} & \textbf{内容} \\
\midrule
\endhead
用例编号：用例名 & UC-01：登记成交 \\
\midrule
执行者 & 销售（主执行者） \\
\midrule
前置条件 & 系统能检测：所选货品已建档入库且处于“在柜”状态（散客可不关联客户档案） \\
\midrule
后置条件 & 新增一条成交记录，必然关联到唯一货品与经手销售；货品状态改为“已售”或“已订” \\
\midrule
涉众利益 & 销售：登记要快、不打断成交；\newline 老板：每笔成交可追溯到货品、客户与时间，周转与营销数据可信；\newline 客户：价格与小票一致，售后有据 \\
\midrule
基本路径 & 1. 销售扫描货品条形码或提供货品编号、成交价、支付方式（全款/定金）与客户（散客可不提供）；\newline 2. 系统校验货品处于“在柜”状态；\newline 3. 系统校验成交价合法（不低于该货品最低售价）；\newline 4. 系统创建成交记录，货品状态改为“已售”（定金成交则标记“已订”）；\newline 5. 系统反馈登记结果 \\
\midrule
扩展 & 2a. 货品编号不存在或状态非“在柜”（已订、已售、预留、已调出、代销中、已送检、待处理）：2a1 系统反馈应重新选择货品；2a2 返回1；\newline 3a. 成交价低于最低售价：3a1 系统要求录入老板授权价及授权说明，无授权时反馈应重新报价；3a2 返回1；\newline 3b. 定金金额超过成交价：3b1 系统反馈应重新输入；3b2 返回1；\newline 4a. 保存失败：4a1 系统反馈应稍后重试；4a2 返回1 \\
\midrule
字段列表 & 货品＝货品编号＋品名＋类别＋料种（自定义）＋产地（自定义）＋器型（自定义）＋产状（山料、籽料、山流水料、戈壁料）＋克重＋尺寸＋证书号＋图片＋来源类型＋来源方＋成本＋最低售价＋入库时间＋状态（在柜、已送检、预留、已售、已订、已调出、代销中、待处理）；\newline 成交记录＝成交编号＋成交时间＋货品＋成交价＋支付方式（全款、定金）＋已收金额＋客户＋经手人 \\
\midrule
业务规则 & ①每件玉器一物一码，一件货品同时只能有一条有效成交记录；\newline ②成交价不得低于该货品最低售价，低于时必须录入老板授权；\newline ③客户明确要求预留但尚未付款的，货品状态置为“预留”，预留期间不得出售给他人；\newline ④客户选择分期留货的，收定金即标记“已订”，货品留店不出库，与预留同样不得再售他人，尾款付清并交付出库后才计为完成销售；\newline ⑤寄卖在同行处的货品状态为“代销中”，货权仍归本店，售出后转为“已售”并与同行结算；\newline ⑥成交时间以系统时间为准，发生时间与录入时间分别记录 \\
\midrule
质量需求 & 可用性：完成一次登记不超过 4 次操作、10 秒内完成；\newline 性能：提交后 1 秒内反馈 \\
\midrule
设计约束 & 柜台端用电脑操作并部署在门店本地网络；供应商、加工师傅与同行商户通过微信小程序访问；货款当面收付（现金或个人收款码），可按需要配小票打印机与条形码标签打印机；不接入在线支付接口 \\
\midrule
\bottomrule
\end{longtable}}
\noindent（2）登记货品入库

{\small
\begin{longtable}{@{}L{2.4cm} L{12.9cm}@{}}
\caption{登记货品入库用例规约}\label{tab:uc-02}\\
\toprule
\textbf{标题} & \textbf{内容} \\
\midrule
\endfirsthead
\toprule
\textbf{标题} & \textbf{内容} \\
\midrule
\endhead
用例编号：用例名 & UC-02：登记货品入库 \\
\midrule
执行者 & 老板（主执行者） \\
\midrule
前置条件 & 来源方（供应商、加工师傅、同行商户）已在基础数据中维护 \\
\midrule
后置条件 & 新增一条货品档案，状态“在柜”，入库时间以系统时间为准 \\
\midrule
涉众利益 & 老板：每件货品都有唯一编号和完整来源、成本信息，周转统计才有依据；\newline 销售：录入成交时按编号即可带出货品资料 \\
\midrule
基本路径 & 1. 老板提供货品信息（品名、类别、料种、产地、器型、克重、尺寸、证书号、图片、来源类型、来源方、成本、最低售价）；\newline 2. 系统校验字段与金额合法；\newline 3. 系统逐件生成货品编号、保存档案并打印条形码标签（印货品编号，系挂于货品）；\newline 4. 系统反馈入库结果与编号清单 \\
\midrule
扩展 & 1a. 来源类型为回收：1a1 系统按回流作价确定成本并标记来源与作价；\newline 1b. 来源类型为调货：1b1 系统标记来源与关联商户，不记成本、只记往来；\newline 2a. 必填字段缺失或金额不合法：2a1 系统反馈应重新输入；2a2 返回1；\newline 3a. 同批次多件（如珠串类）：3a1 系统逐件建档、共用批次信息，编号各不相同 \\
\midrule
字段列表 & 货品来源＝来源类型（采购、回收、调货、加工完工、寄售）＋来源方＋来源单据；\newline 入库记录＝入库时间＋验收人＋货品清单 \\
\midrule
业务规则 & ①一件货品一个编号，重件需拆分为多条档案；\newline ②回收货品按回流作价入成本；\newline ③调货调入不记成本，只记往来；\newline ④加工完工按料加工合计入成本；\newline ⑤寄售货品按寄售价入成本，作为结算依据 \\
\midrule
质量需求 & 可用性：逐件建档不超过 30 秒/件 \\
\midrule
设计约束 & 与现有纸质台账并行运行一个月，双轨核对无误后停用台账 \\
\midrule
\bottomrule
\end{longtable}}
\noindent（3）登记退货

{\small
\begin{longtable}{@{}L{2.4cm} L{12.9cm}@{}}
\caption{登记退货用例规约}\label{tab:uc-03}\\
\toprule
\textbf{标题} & \textbf{内容} \\
\midrule
\endfirsthead
\toprule
\textbf{标题} & \textbf{内容} \\
\midrule
\endhead
用例编号：用例名 & UC-03：登记退货 \\
\midrule
执行者 & 销售（主执行者） \\
\midrule
前置条件 & 待退货的成交记录存在且未标记过退货 \\
\midrule
后置条件 & 原成交记录标记“已退货”，货品状态改回“在柜”，当期销售统计冲减 \\
\midrule
涉众利益 & 客户：真假与质量问题能退、退款及时，不喜欢可协商换货；\newline 老板：退货限于成交后 7 天内且限真假与质量原因，防止随意退货冲减业绩；\newline 销售：不用手工调整台账，换货协商有记录可依 \\
\midrule
基本路径 & 1. 销售提供原成交记录与退货原因；\newline 2. 系统校验该成交存在且未标记过退货，且在成交后 7 天内、退货原因属于可退范围；\newline 3. 系统新增退货记录，货品状态改回“在柜”；\newline 4. 系统反馈登记结果 \\
\midrule
扩展 & 2a. 原成交记录不存在或已退货：2a1 系统反馈不允许重复退货；\newline 2b. 已超过成交后 7 天：2b1 系统反馈已超过退货期限，不予办理；\newline 2c. 退货原因不属于可退范围：2c1 系统反馈该原因不在可退范围内，不予办理；\newline 3a. 退货价与原成交价不一致：3a1 销售说明差额原因，系统将差额计入退货记录 \\
\midrule
字段列表 & 退货记录＝退货时间＋原成交记录＋退货原因（真假存疑、质量问题）＋退款金额＋经手人 \\
\midrule
业务规则 & ①退货必须关联原成交记录，无记录不走系统退货；\newline ②退货期限自成交日起 7 天，超期不予办理；\newline ③仅真假存疑或质量问题可退；圈口不合、主观不满意等原因不予退款，由销售协商换货处理；\newline ④退货后货品状态改回“在柜”，可再次出售；\newline ⑤当月退货冲减当月销售统计 \\
\midrule
质量需求 & 可用性：完成一次退货登记不超过 3 次操作 \\
\midrule
设计约束 & 退款为当面现金退回或按原收款方式退回，不接入在线支付接口 \\
\midrule
\bottomrule
\end{longtable}}
\noindent（4）维护客户信息

{\small
\begin{longtable}{@{}L{2.4cm} L{12.9cm}@{}}
\caption{维护客户信息用例规约}\label{tab:uc-04}\\
\toprule
\textbf{标题} & \textbf{内容} \\
\midrule
\endfirsthead
\toprule
\textbf{标题} & \textbf{内容} \\
\midrule
\endhead
用例编号：用例名 & UC-04：维护客户信息 \\
\midrule
执行者 & 销售（主执行者） \\
\midrule
前置条件 & 新增客户时可提供姓名与联系方式 \\
\midrule
后置条件 & 客户档案新增或更新，客户编号唯一；历史成交记录中的客户引用保持不变 \\
\midrule
涉众利益 & 销售：回头客到店能快速调出资料与购买记录，手机号可一键复制去加微信；\newline 老板：客户资料完整真实，来源渠道与标签能看出哪个渠道有效、谁是核心客户；\newline 客户：个人信息只用于门店经营，不外泄 \\
\midrule
基本路径 & 1. 销售提供客户姓名与联系方式，并选择来源渠道与标签；\newline 2. 系统校验手机号未与其他客户重复；\newline 3. 系统保存客户档案并生成客户编号；\newline 4. 系统反馈维护结果 \\
\midrule
扩展 & 2a. 手机号已存在：2a1 系统带出原客户档案，销售可补充修改；2a2 返回1；\newline 3a. 必填项缺失：3a1 系统反馈应补全后保存；3a2 返回1 \\
\midrule
字段列表 & 客户＝客户编号＋姓名＋手机号＋微信号＋来源渠道（闲鱼、微信朋友圈、朋友介绍、线下、直播、老客户复购、其他）＋标签（大客户、懂行、爽快、要还价）＋偏好（品类、圈口、预算）＋购买记录＋跟进记录＋备注 \\
\midrule
业务规则 & ①同一手机号视为同一客户，不重复建档；\newline ②购买记录由成交登记自动带出，不手工录入；\newline ③来源渠道与标签从固定选项中选择，不自由填写，以便按渠道统计 \\
\midrule
质量需求 & 可用性：按编号、姓名、手机号、微信号、标签任一条件查询，3 秒内反馈；\newline 安全性：手机号等联系方式仅老板与销售可见 \\
\midrule
设计约束 & 不做短信/微信营销推送，仅作查询与筛选；不做会员等级、积分与优惠券；不做客户画像算法打分，核心客户按累计成交与标签筛选得出 \\
\midrule
\bottomrule
\end{longtable}}
\noindent（5）登记送检与证书

{\small
\begin{longtable}{@{}L{2.4cm} L{12.9cm}@{}}
\caption{登记送检与证书用例规约}\label{tab:uc-05}\\
\toprule
\textbf{标题} & \textbf{内容} \\
\midrule
\endfirsthead
\toprule
\textbf{标题} & \textbf{内容} \\
\midrule
\endhead
用例编号：用例名 & UC-05：登记送检与证书 \\
\midrule
执行者 & 销售（主执行者） \\
\midrule
前置条件 & 货品已建档入库且状态为“在柜”或“预留” \\
\midrule
后置条件 & 新增一条送检记录，货品状态改为“已送检”；证书信息保存后与货品关联，货品状态恢复送检前的状态（在柜或预留） \\
\midrule
涉众利益 & 销售：证书编号与货品一一对应，客户索证时随时可调，销售时直接出示证书；\newline 老板：不用亲自跑送检，只需安排与把关；\newline 鉴定机构：送检信息与机构回执一致 \\
\midrule
基本路径 & 1. 销售提供货品编号、送检机构与送检日期；\newline 2. 系统校验货品状态为“在柜”或“预留”；\newline 3. 系统新增送检记录，货品状态改为“已送检”；\newline 4. 销售录入鉴定机构返回的证书编号与检测结论；\newline 5. 系统保存证书信息并关联货品，货品状态恢复送检前状态；\newline 6. 系统反馈登记结果 \\
\midrule
扩展 & 2a. 货品状态非“在柜”且非“预留”（已售、已调出、代销中）：2a1 系统反馈应重新选择货品；2a2 返回1；\newline 4a. 检测结论为不合格：4a1 系统保存结论并把货品状态改为“待处理”，由老板决定退回供应商或停售 \\
\midrule
字段列表 & 送检记录＝送检时间＋货品＋送检机构＋送检人；\newline 证书＝证书编号＋鉴定机构＋检测结论＋签发日期＋证书二维码（机构印制，扫码可查线上鉴定证书） \\
\midrule
业务规则 & ①证书编号与货品一一对应，不得复用；\newline ②送检期间货品状态为“已送检”，不得出售；\newline ③证书编号录入后不可修改，只能作废重录；\newline ④鉴定证书只判真假，产地由老板凭眼力判断并记录在货品档案 \\
\midrule
质量需求 & 性能：提交后 1 秒内反馈 \\
\midrule
设计约束 & 鉴定机构信息从基础数据带出，不与机构系统对接 \\
\midrule
\bottomrule
\end{longtable}}
\noindent（6）查看销售与周转统计

{\small
\begin{longtable}{@{}L{2.4cm} L{12.9cm}@{}}
\caption{查看销售与周转统计用例规约}\label{tab:uc-06}\\
\toprule
\textbf{标题} & \textbf{内容} \\
\midrule
\endfirsthead
\toprule
\textbf{标题} & \textbf{内容} \\
\midrule
\endhead
用例编号：用例名 & UC-06：查看销售与周转统计 \\
\midrule
执行者 & 老板（主执行者） \\
\midrule
前置条件 & 系统内已有货品入库记录与成交记录 \\
\midrule
后置条件 & 只读查询，不改变任何数据 \\
\midrule
涉众利益 & 老板：统计数字能与纸质台账核对一致，按月对比销售额与周转周期两个愿景指标；\newline 销售：可查自己经手的成交，数据不被他人改动 \\
\midrule
基本路径 & 1. 老板选择统计时段（月/季/年）；\newline 2. 系统汇总该时段销售额、销售客户数与货品加权周转周期；\newline 3. 系统按品种与价位筛选周转最慢的货品；\newline 4. 系统按来源渠道汇总成交，并列出累计贡献靠前的客户；\newline 5. 系统反馈统计结果 \\
\midrule
扩展 & 2a. 时段内无记录：2a1 系统反馈该时段无数据；2a2 返回1 \\
\midrule
字段列表 & 统计口径＝月度销售额＋销售客户数（去重）＋货品加权周转周期（按货品总金额加权）＋来源渠道成交分布＋核心客户贡献（累计成交额与占比） \\
\midrule
业务规则 & ①销售额按成交时间归属月份，退货冲减；\newline ②周转周期从货品入库到销售完成并收回货款，按货品总金额加权 \\
\midrule
质量需求 & 性能：统计查询 2 秒内反馈 \\
\midrule
设计约束 & 统计界面只读，不提供数据导出 \\
\midrule
\bottomrule
\end{longtable}}
\noindent（7）维护基础数据

{\small
\begin{longtable}{@{}L{2.4cm} L{12.9cm}@{}}
\caption{维护基础数据用例规约}\label{tab:uc-07}\\
\toprule
\textbf{标题} & \textbf{内容} \\
\midrule
\endfirsthead
\toprule
\textbf{标题} & \textbf{内容} \\
\midrule
\endhead
用例编号：用例名 & UC-07：维护基础数据 \\
\midrule
执行者 & 老板（主执行者） \\
\midrule
前置条件 & 无（身份验证由 UC-08 提供） \\
\midrule
后置条件 & 来源方（供应商/加工师傅/同行商户）、鉴定机构与货品类别等参照数据新增或更新，供入库与成交登记选择 \\
\midrule
涉众利益 & 老板：参照数据齐全，录入时直接选择、不用每次手输；\newline 销售：选择项稳定统一，统计不因写法不同出现偏差 \\
\midrule
基本路径 & 1. 老板选择要维护的数据类型（来源方/鉴定机构/货品类别，来源方以类别区分供应商/加工师傅/同行商户）；\newline 2. 老板提供新增、修改或删除的内容；\newline 3. 系统校验名称在同一类别内未重复；\newline 4. 系统保存并反馈维护结果 \\
\midrule
扩展 & 3a. 名称重复：3a1 系统反馈应更换名称；3a2 返回1；\newline 4a. 删除已被引用的数据项：4a1 系统反馈该数据项已被成交或入库记录引用，只能停用；4a2 返回1 \\
\midrule
字段列表 & 基础数据＝来源方（类别＋名称＋联系人＋联系电话）＋鉴定机构＋货品类别 \\
\midrule
业务规则 & ①基础数据被引用后不可物理删除，只能停用；\newline ②供应商类型区分玉料/成品/回流 \\
\midrule
质量需求 & 可用性：新增一条基础数据不超过 3 次操作 \\
\midrule
设计约束 & 基础数据仅老板可维护，销售只读 \\
\midrule
\bottomrule
\end{longtable}}
\noindent（8）登录

{\small
\begin{longtable}{@{}L{2.4cm} L{12.9cm}@{}}
\caption{登录用例规约}\label{tab:uc-08}\\
\toprule
\textbf{标题} & \textbf{内容} \\
\midrule
\endfirsthead
\toprule
\textbf{标题} & \textbf{内容} \\
\midrule
\endhead
用例编号：用例名 & UC-08：登录 \\
\midrule
执行者 & 老板、销售、合作商户、供应商、加工师傅（被 UC-01～UC-07 与 UC-11、UC-12、UC-14～UC-26、UC-28 以 \code{<<include>>} 包含的公共用例；客户无需登录） \\
\midrule
前置条件 & 用户账号已由老板建立 \\
\midrule
后置条件 & 登录成功后系统按角色开放对应功能；后续关键操作归属到该用户名下 \\
\midrule
涉众利益 & 老板：职责分离，销售无法使用授权价、基础数据维护等老板功能；\newline 销售：口令自行保管，经手记录归属清晰 \\
\midrule
基本路径 & 1. 执行者提供用户名与口令；\newline 2. 系统校验用户名存在且口令正确；\newline 3. 系统按角色开放对应功能并反馈登录结果 \\
\midrule
扩展 & 2a. 用户名不存在或口令错误：2a1 系统反馈应重试；2a2 返回1；\newline 2b. 同一账号连续 5 次口令错误：2b1 系统锁定账号 10 分钟，并提示联系老板解锁 \\
\midrule
字段列表 & 用户＝用户名＋姓名＋口令（加密保存）＋角色（老板、销售、合作商户、供应商、加工师傅）＋状态（启用、停用） \\
\midrule
业务规则 & ①连续 5 次口令错误锁定 10 分钟；\newline ②角色决定登录后的功能范围 \\
\midrule
质量需求 & 安全性：口令加密保存；\newline 性能：登录验证 1 秒内反馈 \\
\midrule
设计约束 & 不做自助注册，账号全部由老板开通 \\
\midrule
\bottomrule
\end{longtable}}
\noindent（9）浏览货品

{\small
\begin{longtable}{@{}L{2.4cm} L{12.9cm}@{}}
\caption{浏览货品用例规约}\label{tab:uc-09}\\
\toprule
\textbf{标题} & \textbf{内容} \\
\midrule
\endfirsthead
\toprule
\textbf{标题} & \textbf{内容} \\
\midrule
\endhead
用例编号：用例名 & UC-09：浏览货品 \\
\midrule
执行者 & 客户（无需登录，也不提供注册） \\
\midrule
前置条件 & 系统内存在状态为“在柜”且允许展示的货品档案 \\
\midrule
后置条件 & 只读查询，不改变任何数据 \\
\midrule
涉众利益 & 客户：随时了解门店在售货品，到店前有预期；\newline 老板：展示范围可控，价格信息不外露；\newline 销售：预约到店的客户带着目标来，接待更省力 \\
\midrule
基本路径 & 1. 客户请求浏览在柜货品；\newline 2. 系统检索状态为“在柜”且允许展示的货品图文信息；\newline 3. 客户选择货品查看详情；\newline 4. 系统反馈货品展示信息（不含任何价格） \\
\midrule
扩展 & 2a. 货品已售出、已调出或已送检：2a1 系统不再展示该货品 \\
\midrule
字段列表 & 展示项＝图片＋品名＋类别＋尺寸/圈口＋描述（不含价格） \\
\midrule
业务规则 & ①仅展示“在柜”货品；\newline ②一物一码对应一页展示 \\
\midrule
质量需求 & 性能：图片加载 2 秒内完成；\newline 可用性：微信小程序或手机浏览器可直接打开 \\
\midrule
设计约束 & 前端不显示价格与成本 \\
\midrule
\bottomrule
\end{longtable}}
\noindent（10）预约查看货品

{\small
\begin{longtable}{@{}L{2.4cm} L{12.9cm}@{}}
\caption{预约查看货品用例规约}\label{tab:uc-10}\\
\toprule
\textbf{标题} & \textbf{内容} \\
\midrule
\endfirsthead
\toprule
\textbf{标题} & \textbf{内容} \\
\midrule
\endhead
用例编号：用例名 & UC-10：预约查看货品 \\
\midrule
执行者 & 客户（无需登录） \\
\midrule
前置条件 & 目标货品处于“在柜”且在展示范围内 \\
\midrule
后置条件 & 新增一条预约记录，关联货品、联系方式与期望到店时段；货品状态不变；预约编号与时段在页面上回显 \\
\midrule
涉众利益 & 客户：到店前确认货品在店、有人接待；\newline 销售：提前知道接待对象与看点；\newline 老板：预约数量是前端引流效果的直接指标 \\
\midrule
基本路径 & 1. 客户选择货品并请求预约查看；\newline 2. 客户提供姓名、手机号与期望到店时段；\newline 3. 系统校验手机号与时段合法，保存预约记录并生成预约编号；\newline 4. 系统反馈预约编号、期望时段与门店地址，并提示客户截图保存 \\
\midrule
扩展 & 2a. 手机号格式不合法：2a1 系统反馈应重新输入；2a2 返回2；\newline 3a. 同一手机号当日重复预约同一货品：3a1 系统提示已有预约，可查看原预约 \\
\midrule
字段列表 & 预约＝预约编号＋货品＋姓名＋手机号＋期望时段＋状态 \\
\midrule
业务规则 & ①同一货品同一时段允许多组预约；\newline ②预约信息仅用于到店接待 \\
\midrule
质量需求 & 性能：提交后 1 秒内反馈 \\
\midrule
设计约束 & 无需登录，手机号作为联系标识 \\
\midrule
\bottomrule
\end{longtable}}
\noindent（11）查看调货价格

{\small
\begin{longtable}{@{}L{2.4cm} L{12.9cm}@{}}
\caption{查看调货价格用例规约}\label{tab:uc-11}\\
\toprule
\textbf{标题} & \textbf{内容} \\
\midrule
\endfirsthead
\toprule
\textbf{标题} & \textbf{内容} \\
\midrule
\endhead
用例编号：用例名 & UC-11：查看调货价格 \\
\midrule
执行者 & 合作商户（主执行者，账号由老板逐个开通） \\
\midrule
前置条件 & 无（身份验证由 UC-08 提供） \\
\midrule
后置条件 & 只读查询，不改变任何数据 \\
\midrule
涉众利益 & 合作商户：自助掌握可调货品与调货价格，减少往返问价；\newline 老板：调货价格只对开通账户的同行可见，价格体系可控；\newline 客户：不受影响，前端不展示价格 \\
\midrule
基本路径 & 1. 合作商户选择货品类别或输入品名关键字；\newline 2. 系统检索状态为“在柜”的货品与对应调货价格；\newline 3. 系统反馈查询结果 \\
\midrule
扩展 & 2a. 货品未开放展示或已调出：2a1 系统不反馈该货品 \\
\midrule
字段列表 & 调货价目＝货品＋类别＋圈口/尺寸＋调货价 \\
\midrule
业务规则 & ①调货价仅对“合作商户”角色可见；\newline ②调货价与零售价分离维护 \\
\midrule
质量需求 & 性能：查询 2 秒内反馈 \\
\midrule
设计约束 & 价目页只读，不提供下载 \\
\midrule
\bottomrule
\end{longtable}}
\noindent（12）开通外部账户

{\small
\begin{longtable}{@{}L{2.4cm} L{12.9cm}@{}}
\caption{开通外部账户用例规约}\label{tab:uc-12}\\
\toprule
\textbf{标题} & \textbf{内容} \\
\midrule
\endfirsthead
\toprule
\textbf{标题} & \textbf{内容} \\
\midrule
\endhead
用例编号：用例名 & UC-12：开通外部账户 \\
\midrule
执行者 & 老板（主执行者） \\
\midrule
前置条件 & 开通对象已在基础数据中维护（来源方以类别区分供应商、加工师傅等） \\
\midrule
后置条件 & 新增一个外部账号并分配角色（可多选），可登录办理各自事项；停用后无法登录 \\
\midrule
涉众利益 & 老板：给谁开账号自己决定，价格与往来信息可控；\newline 合作商户：拿到账号即可自助查价，不用每次电话问价；\newline 供应商与加工师傅：线上提交申请与预约，减少往返沟通 \\
\midrule
基本路径 & 1. 老板选择要开通的角色（可多选）并提供账号信息（用户名与初始口令）；\newline 2. 系统校验用户名未重复；\newline 3. 系统创建账号并分配所选角色；\newline 4. 系统反馈开通结果 \\
\midrule
扩展 & 2a. 用户名已存在：2a1 系统反馈应更换用户名；2a2 返回1 \\
\midrule
字段列表 & 用户＝用户名＋姓名＋口令（加密保存）＋状态（启用、停用）；\newline 用户-角色＝用户＋角色（合作商户、供应商、加工师傅等，可多选） \\
\midrule
业务规则 & ①外部账号由老板逐个开通与停用，不能自行注册；\newline ②合作商户角色只有查看调货价格的权限，供应商与加工师傅角色只有各自协同事项的权限，均无内部管理功能；\newline ③口令加密保存，首次登录可修改 \\
\midrule
质量需求 & 安全性：口令加密保存；\newline 可用性：开通一个账号不超过 3 次操作 \\
\midrule
设计约束 & 复用 UC-08 的登录入口，角色决定登录后的功能范围 \\
\midrule
\bottomrule
\end{longtable}}
\noindent（13）预定商品

{\small
\begin{longtable}{@{}L{2.4cm} L{12.9cm}@{}}
\caption{预定商品用例规约}\label{tab:uc-13}\\
\toprule
\textbf{标题} & \textbf{内容} \\
\midrule
\endfirsthead
\toprule
\textbf{标题} & \textbf{内容} \\
\midrule
\endhead
用例编号：用例名 & UC-13：预定商品 \\
\midrule
执行者 & 客户（无需登录） \\
\midrule
前置条件 & 目标货品处于“在柜”且在展示范围内，且当前没有其他有效预定 \\
\midrule
后置条件 & 新增一条预定记录，关联货品与联系方式；货品状态不变，不锁货 \\
\midrule
涉众利益 & 客户：先把看中的货留住，到店有货可看；\newline 销售：知道哪件货被预定，接待有准备；\newline 老板：预定数据反映真实到店意向，防止恶意占货 \\
\midrule
基本路径 & 1. 客户选择货品并请求预定；\newline 2. 客户提供姓名与手机号；\newline 3. 系统校验该手机号的预定额度未满，且该货品没有其他有效预定；\newline 4. 系统保存预定记录并生成预定编号；\newline 5. 系统反馈预定结果 \\
\midrule
扩展 & 3a. 该手机号预定额度已满（占用 3 条）：3a1 系统反馈预定额度已满，暂不可预定；3a2 返回2；\newline 3b. 该货品已有其他有效预定：3b1 系统反馈货品已被预定；3b2 返回2 \\
\midrule
字段列表 & 预定记录＝预定编号＋货品＋姓名＋手机号＋提交时间＋状态（待到店、已联系、已达成、已爽约、已取消） \\
\midrule
业务规则 & ①一个手机号最多同时占用 3 条预定额度；\newline ②预定如约达成（按时到店，由销售确认）后释放 1 条额度；\newline ③爽约或客户取消不释放额度；\newline ④同一货品同时只能有 1 条有效预定；\newline ⑤预定超过约定到店时间未到店（含已联系但未到店），系统标记爽约，该条额度不释放；\newline ⑥客户预定后联系销售并确定到店的，销售将该货品置为“预留”；到店看过暂不决定的，销售取消预留，货品回“在柜” \\
\midrule
质量需求 & 性能：提交后 1 秒内反馈；\newline 可用性：完成一次预定不超过 3 次操作 \\
\midrule
设计约束 & 预定入口从货品展示页进入，剩余预定额度在页面上可见 \\
\midrule
\bottomrule
\end{longtable}}
\noindent（14）提交货品出让申请

{\small
\begin{longtable}{@{}L{2.4cm} L{12.9cm}@{}}
\caption{提交货品出让申请用例规约}\label{tab:uc-14}\\
\toprule
\textbf{标题} & \textbf{内容} \\
\midrule
\endfirsthead
\toprule
\textbf{标题} & \textbf{内容} \\
\midrule
\endhead
用例编号：用例名 & UC-14：提交货品出让申请 \\
\midrule
执行者 & 供应商（含玉料供应商、成品供应商、回流出让方，账号由老板开通） \\
\midrule
前置条件 & 无（身份验证由 UC-08 提供） \\
\midrule
后置条件 & 新增一条出让申请，状态“待审核”，等待老板审核 \\
\midrule
涉众利益 & 供应商：在线提交货品信息，不用上门先谈一轮；\newline 老板：申请集中在线上，有空再审核，货源线索不散失；\newline 销售：审核通过的货源到店后按常规流程接待 \\
\midrule
基本路径 & 1. 供应商选择出让类型（玉料/成品/回流）并填写货品信息（品名、数量、期望出让价、描述），提供料子图片；\newline 2. 系统校验必填字段完整；\newline 3. 系统保存申请并生成申请编号，状态置为“待审核”；\newline 4. 系统反馈提交结果 \\
\midrule
扩展 & 2a. 必填字段缺失：2a1 系统反馈应补全后提交；2a2 返回1 \\
\midrule
字段列表 & 出让申请＝申请编号＋供应商＋出让类型（玉料、成品、回流）＋品名＋数量＋期望出让价＋描述＋料子图片＋提交时间＋状态（待审核、已通过、已驳回） \\
\midrule
业务规则 & ①玉料供应商、成品供应商与回流出让方共用同一申请入口，类型取自供应商档案；\newline ②申请经老板审核通过后，货品实际到店时按 UC-02 建档入库；\newline ③期望出让价仅供审核参考，不对外展示 \\
\midrule
质量需求 & 性能：提交后 1 秒内反馈；\newline 可用性：完成一次出让申请不超过 5 次操作 \\
\midrule
设计约束 & 申请入口在供应商登录后的界面，与调货价目查询互不可见 \\
\midrule
\bottomrule
\end{longtable}}
\noindent（15）预约上门时间

{\small
\begin{longtable}{@{}L{2.4cm} L{12.9cm}@{}}
\caption{预约上门时间用例规约}\label{tab:uc-15}\\
\toprule
\textbf{标题} & \textbf{内容} \\
\midrule
\endfirsthead
\toprule
\textbf{标题} & \textbf{内容} \\
\midrule
\endhead
用例编号：用例名 & UC-15：预约上门时间 \\
\midrule
执行者 & 供应商、加工师傅（账号由老板开通） \\
\midrule
前置条件 & 供应商有已通过的出让申请，或加工师傅有已确认的完工交付申请 \\
\midrule
后置条件 & 新增一条上门预约，关联预约方、事项与时段 \\
\midrule
涉众利益 & 供应商与加工师傅：约好时间再来，不白跑；\newline 老板：交接时间集中安排，减少临时接待；\newline 销售：知道什么时间有什么交接，提前准备 \\
\midrule
基本路径 & 1. 预约方选择要办理的事项（出让交货/完工交付）；\newline 2. 预约方提供期望上门时段；\newline 3. 系统校验该时段没有其他上门预约冲突，保存预约；\newline 4. 系统反馈预约结果 \\
\midrule
扩展 & 3a. 时段已有其他上门预约：3a1 系统反馈应更换时段；3a2 返回2 \\
\midrule
字段列表 & 上门预约＝预约编号＋预约方＋事项＋期望时段＋提交时间＋状态（待上门、已完成、已取消） \\
\midrule
业务规则 & ①预约事项必须关联一条已通过的出让申请或已确认的完工交付申请；\newline ②同一时段只安排一次上门交接；\newline ③预约由老板在系统内确认，改期需重新预约 \\
\midrule
质量需求 & 性能：提交后 1 秒内反馈；\newline 可用性：预约不超过 3 次操作 \\
\midrule
设计约束 & 上门交接与客户到店预约统一安排时段，避免相互冲突 \\
\midrule
\bottomrule
\end{longtable}}
\noindent（16）查看加工任务

{\small
\begin{longtable}{@{}L{2.4cm} L{12.9cm}@{}}
\caption{查看加工任务用例规约}\label{tab:uc-16}\\
\toprule
\textbf{标题} & \textbf{内容} \\
\midrule
\endfirsthead
\toprule
\textbf{标题} & \textbf{内容} \\
\midrule
\endhead
用例编号：用例名 & UC-16：查看加工任务 \\
\midrule
执行者 & 加工师傅（账号由老板开通） \\
\midrule
前置条件 & 无（身份验证由 UC-08 提供） \\
\midrule
后置条件 & 只读查询，不改变任何数据 \\
\midrule
涉众利益 & 加工师傅：随时掌握待加工与在加工任务、要求与工费；\newline 老板：不用电话逐个通知师傅；\newline 销售：加工进度可查，客户问货时能答复 \\
\midrule
基本路径 & 1. 加工师傅请求查看加工任务；\newline 2. 系统检索分派给该师傅的任务（货品、加工要求、工费、期限、状态）；\newline 3. 系统反馈查询结果 \\
\midrule
扩展 & 2a. 没有任何任务：2a1 系统反馈当前无加工任务 \\
\midrule
字段列表 & 加工任务＝任务编号＋货品＋加工要求＋工费＋期限＋状态（待加工、在加工、待交付、已交付） \\
\midrule
业务规则 & ①任务由老板在安排加工时生成，加工师傅只能查看，不能改任务内容；\newline ②工费金额对师傅可见，结算仍由老板在系统内记录 \\
\midrule
质量需求 & 性能：查询 2 秒内反馈；\newline 可靠性：任务状态与老板端记录一致 \\
\midrule
设计约束 & 加工师傅与内部账号共用登录入口，登录后只见任务相关功能 \\
\midrule
\bottomrule
\end{longtable}}
\noindent（17）提交完工交付申请

{\small
\begin{longtable}{@{}L{2.4cm} L{12.9cm}@{}}
\caption{提交完工交付申请用例规约}\label{tab:uc-17}\\
\toprule
\textbf{标题} & \textbf{内容} \\
\midrule
\endfirsthead
\toprule
\textbf{标题} & \textbf{内容} \\
\midrule
\endhead
用例编号：用例名 & UC-17：提交完工交付申请 \\
\midrule
执行者 & 加工师傅（账号由老板开通） \\
\midrule
前置条件 & 存在状态为“在加工”的本人任务 \\
\midrule
后置条件 & 任务状态改为“待交付”，等待老板确认并安排入库 \\
\midrule
涉众利益 & 加工师傅：完工即在线报备，不用等老板有空再交接；\newline 老板：完工信息第一时间可见，及时安排交接入库；\newline 销售：货品加工完成的信息可查，便于对客户答复 \\
\midrule
基本路径 & 1. 加工师傅选择已完工的加工任务；\newline 2. 加工师傅提供加工前、加工中与完工的实物图与完工说明；\newline 3. 系统校验任务状态为“在加工”，更新为“待交付”；\newline 4. 系统反馈提交结果 \\
\midrule
扩展 & 3a. 任务状态不是“在加工”：3a1 系统反馈不允许重复提交；3a2 返回1 \\
\midrule
字段列表 & 完工交付申请＝申请编号＋任务编号＋完工说明＋过程图（加工前/加工中/完工）＋提交时间＋状态（待确认、已确认） \\
\midrule
业务规则 & ①完工交付申请提交后，货品入库仍由老板按 UC-02 建档，加工师傅不录入货品档案；\newline ②老板确认完工后，师傅通过 UC-15 预约上门时间交货；\newline ③一条任务只能有一条待确认的完工交付申请 \\
\midrule
质量需求 & 性能：提交后 1 秒内反馈；\newline 可靠性：完工状态与老板端一致 \\
\midrule
设计约束 & 提交入口从加工任务列表进入；老板确认后按 UC-15 预约交接时段 \\
\midrule
\bottomrule
\end{longtable}}
\noindent（18）审核出让申请

{\small
\begin{longtable}{@{}L{2.4cm} L{12.9cm}@{}}
\caption{审核出让申请用例规约}\label{tab:uc-18}\\
\toprule
\textbf{标题} & \textbf{内容} \\
\midrule
\endfirsthead
\toprule
\textbf{标题} & \textbf{内容} \\
\midrule
\endhead
用例编号：用例名 & UC-18：审核出让申请 \\
\midrule
执行者 & 老板（主执行者） \\
\midrule
前置条件 & 存在状态为“待审核”的出让申请 \\
\midrule
后置条件 & 申请状态改为“已通过”或“已驳回”；通过的申请在货品到店时转入入库流程 \\
\midrule
涉众利益 & 老板：货源线索集中审核，通过与否自己把关；\newline 供应商：审核结果在线可查，不用反复电话确认；\newline 销售：通过的货源到店后按常规流程接待 \\
\midrule
基本路径 & 1. 老板请求查看待审核的出让申请；\newline 2. 系统反馈申请详情与期望出让价；\newline 3. 老板选择通过或驳回（驳回须填写原因）；\newline 4. 系统保存审核结果并反馈 \\
\midrule
扩展 & 3a. 驳回未填原因：3a1 系统反馈应填写原因后提交；3a2 返回3 \\
\midrule
字段列表 & 出让申请（含审核）＝出让申请字段＋审核意见＋审核时间＋审核人 \\
\midrule
业务规则 & ①审核通过的申请，货品实际到店时按 UC-02 建档，来源类型与出让类型一致；\newline ②驳回的申请保留记录，不删除；\newline ③审核只能由老板执行，销售可查看但不能审核 \\
\midrule
质量需求 & 性能：审核提交后 1 秒内反馈 \\
\midrule
设计约束 & 审核入口在老板登录后的界面，待审核列表按提交时间排序；通过后由供应商按 UC-15 预约上门交接 \\
\midrule
\bottomrule
\end{longtable}}

\noindent（19）登记货款记录

{\small
\begin{longtable}{@{}L{2.4cm} L{12.9cm}@{}}
\caption{登记货款记录用例规约}\label{tab:uc-19}\\
\toprule
\textbf{标题} & \textbf{内容} \\
\midrule
\endfirsthead
\toprule
\textbf{标题} & \textbf{内容} \\
\midrule
\endhead
用例编号：用例名 & UC-19：登记货款记录 \\
\midrule
执行者 & 老板（主执行者） \\
\midrule
前置条件 & 存在对应业务单据（采购入库、调货往来或加工任务） \\
\midrule
后置条件 & 新增一条货款记录（业务类型、关联单据、金额、收付方向、经手人），货款往来可追溯 \\
\midrule
涉众利益 & 老板：每笔货款可追溯到来源单据与经手人，资金往来一本账；\newline
供应商与加工师傅：货款收付有据可查，结算不扯皮；\newline
销售：不承担记账职责 \\
\midrule
基本路径 & 1. 老板选择业务类型（采购货款/调货货款/加工工费/寄售货款）并关联对应单据；\newline
2. 老板录入金额与收付方向（支付/收取）；\newline
3. 系统校验金额与单据一致；\newline
4. 系统保存货款记录；\newline
5. 系统反馈登记结果 \\
\midrule
扩展 & 3a. 金额与单据不一致：3a1 系统反馈应核对后重新录入；3a2 返回2 \\
\midrule
字段列表 & 货款记录＝记录编号＋业务类型（采购货款、调货货款、加工工费、寄售货款）＋关联单据＋金额＋收付方向（支付、收取）＋经手人＋时间 \\
\midrule
业务规则 & ①每条货款记录必须关联来源单据；\newline ②收付方向区分支付与收取；\newline ③记录只增不改，登记错误时作废重录 \\
\midrule
质量需求 & 可用性：登记一条不超过 3 次操作 \\
\midrule
设计约束 & 与统计口径衔接，付款计入成本侧、收款计入销售侧 \\
\midrule
\bottomrule
\end{longtable}}

\noindent（20）维护调货记录

{\small
\begin{longtable}{@{}L{2.4cm} L{12.9cm}@{}}
\caption{维护调货记录用例规约}\label{tab:uc-20}\\
\toprule
\textbf{标题} & \textbf{内容} \\
\midrule
\endfirsthead
\toprule
\textbf{标题} & \textbf{内容} \\
\midrule
\endhead
用例编号：用例名 & UC-20：维护调货记录 \\
\midrule
执行者 & 老板（主执行者） \\
\midrule
前置条件 & 存在调货业务（合作商户调入或调出） \\
\midrule
后置条件 & 新增或更新调货往来记录与调货库存记录，调货货品去向可追溯 \\
\midrule
涉众利益 & 老板：调货往来与库存一本账，结算不扯皮；\newline
合作商户：往来记录准确，对账有依据；\newline
销售：库存状态及时更新，接待有准备 \\
\midrule
基本路径 & 1. 老板选择调货业务（调入/调出）并确认货品与数量；\newline
2. 系统新增调货往来记录；\newline
3. 老板登记货品去向（售出/退回），系统更新调货库存记录；\newline
4. 系统反馈登记结果 \\
\midrule
扩展 & 3a. 货品状态与调货记录不符：3a1 系统反馈应核对后重新登记；3a2 返回3 \\
\midrule
字段列表 & 调货记录＝记录编号＋方向（调入、调出）＋合作商户＋货品＋数量＋状态（在调、已售、已退回）＋时间 \\
\midrule
业务规则 & ①往来记录与库存记录联动更新；\newline ②调给同行的货以客户鉴赏期（3 天或 7 天）届满为结账时点；\newline ③鉴赏期内客户确定留货即按已售出处理，触发与商户的货款结算 \\
\midrule
质量需求 & 可靠性：往来记录与老板端一致 \\
\midrule
设计约束 & 调货记录仅老板可维护，销售只读 \\
\midrule
\bottomrule
\end{longtable}}

\noindent（21）登记寄售合同

{\small
\begin{longtable}{@{}L{2.4cm} L{12.9cm}@{}}
\caption{登记寄售合同用例规约}\label{tab:uc-21}\\
\toprule
\textbf{标题} & \textbf{内容} \\
\midrule
\endfirsthead
\toprule
\textbf{标题} & \textbf{内容} \\
\midrule
\endhead
用例编号：用例名 & UC-21：登记寄售合同 \\
\midrule
执行者 & 老板（主执行者） \\
\midrule
前置条件 & 供应商在基础数据中已建档 \\
\midrule
后置条件 & 新增一份状态为“生效”的寄售合同与寄售明细，货品以寄售方式入库（来源类型为寄售、成本为寄售价、状态为在柜） \\
\midrule
涉众利益 & 老板：寄售货品有合同与寄售价为据，售出后结算不扯皮；\newline
供应商：货品在店期间去向可查，售出按寄售价回款；\newline
销售：寄售货品与自有货品同场展示，无须区分维护 \\
\midrule
基本路径 & 1. 老板选择供应商并填写合同信息（到期日、结算方式）；\newline
2. 系统校验供应商存在，创建寄售合同；\newline
3. 老板逐件录入寄售货品（品名、类别、克重与寄售价）；\newline
4. 系统按寄售方式建档入库并创建寄售明细；\newline
5. 系统写入库存流水；\newline
6. 系统反馈登记结果（含合同编号） \\
\midrule
扩展 & 2a. 供应商不存在：2a1 系统反馈先在维护基础数据中建档；2a2 返回1 \\
\midrule
字段列表 & 寄售合同＝合同编号＋供应商＋签订时间＋到期日＋结算方式（售出即结、按月结算）＋状态（生效、已终止、已到期）；\newline 寄售明细＝合同＋货品＋寄售价＋结算状态（未结算、已结算） \\
\midrule
业务规则 & ①一份合同可含多件寄售货品，一件货品同一时间只挂一份未结算的寄售明细；\newline ②寄售货品的成本记为寄售价，作为结算依据；\newline ③合同终止或到期后，未结算明细对应的货品退回供应商 \\
\midrule
质量需求 & 可用性：登记一件寄售货品不超过 4 次操作 \\
\midrule
设计约束 & 寄售货品与自有货品在柜面同场展示，售出流程与 UC-01 相同 \\
\midrule
\bottomrule
\end{longtable}}

\noindent（22）结算寄售货款

{\small
\begin{longtable}{@{}L{2.4cm} L{12.9cm}@{}}
\caption{结算寄售货款用例规约}\label{tab:uc-22}\\
\toprule
\textbf{标题} & \textbf{内容} \\
\midrule
\endfirsthead
\toprule
\textbf{标题} & \textbf{内容} \\
\midrule
\endhead
用例编号：用例名 & UC-22：结算寄售货款 \\
\midrule
执行者 & 老板（主执行者） \\
\midrule
前置条件 & 存在已售出且未结算的寄售明细 \\
\midrule
后置条件 & 新增一条台账结算（类型为寄售），对应寄售明细置为已结算；寄售合同状态同步更新 \\
\midrule
涉众利益 & 老板：按寄售价与供应商结清，账实一致；\newline
供应商：回款有结算单为据；\newline
销售：不承担结算职责 \\
\midrule
基本路径 & 1. 老板选择已售出的寄售货品；\newline
2. 系统校验该明细未结算；\newline
3. 系统按寄售价创建台账结算（往来方为供应商）；\newline
4. 系统将寄售明细置为已结算；\newline
5. 系统按结算或退回结果更新寄售合同状态（明细全部结清置已结清，到期未售出退回的置已到期）；\newline
6. 系统反馈结算结果（含结算编号） \\
\midrule
扩展 & 2a. 明细已结算：2a1 系统反馈该货品已结清；2a2 返回1 \\
\midrule
字段列表 & 台账结算＝结算编号＋往来方＋结算类型（调货、寄售）＋结算金额＋结算方式（现金、转账、冲抵）＋结算时间＋经手人 \\
\midrule
业务规则 & ①结算金额为寄售价；\newline ②一张结算单只对应一张业务单据（寄售明细或调货记录）；\newline ③结算完成后资金流水在 UC-19 货款记录登记 \\
\midrule
质量需求 & 可用性：结算一笔不超过 3 次操作 \\
\midrule
设计约束 & 寄售与调货结算共用一本台账 \\
\midrule
\bottomrule
\end{longtable}}

\noindent（23）结算调货台账

{\small
\begin{longtable}{@{}L{2.4cm} L{12.9cm}@{}}
\caption{结算调货台账用例规约}\label{tab:uc-23}\\
\toprule
\textbf{标题} & \textbf{内容} \\
\midrule
\endfirsthead
\toprule
\textbf{标题} & \textbf{内容} \\
\midrule
\endhead
用例编号：用例名 & UC-23：结算调货台账 \\
\midrule
执行者 & 老板（主执行者） \\
\midrule
前置条件 & 存在未结算的调货记录 \\
\midrule
后置条件 & 按所选调货记录逐条新增台账结算（类型为调货，一单一据），所选记录均置为已结算 \\
\midrule
涉众利益 & 老板：按商户对账，鉴赏期届满即结清；\newline
合作商户：往来与结算有据可查，不扯皮；\newline
销售：不承担对账职责 \\
\midrule
基本路径 & 1. 老板选择合作商户，系统列出未结算调货记录；\newline
2. 老板勾选本次结算的记录并确认结算方式；\newline
3. 系统按所选记录逐条创建台账结算，每条金额取对应调货记录，结算页面显示金额合计；\newline
4. 系统将所选调货记录置为已结算；\newline
5. 系统反馈结算结果（含各条结算编号） \\
\midrule
扩展 & 1a. 无未结算记录：1a1 系统反馈该商户往来已结清；1a2 用例结束 \\
\midrule
字段列表 & 同台账结算（结算类型为调货），每条结算金额取对应调货记录，结算页面显示所选记录金额合计 \\
\midrule
业务规则 & ①每条结算金额取对应调货记录金额，结算页面显示所选记录金额合计；\newline ②已结算的调货记录不再出现在结算清单；\newline ③与 UC-20 规则②③衔接，鉴赏期届满为结账时点 \\
\midrule
质量需求 & 可靠性：结算后单据与商户端记录一致 \\
\midrule
设计约束 & 结算清单支持一次多笔 \\
\midrule
\bottomrule
\end{longtable}}

\noindent（24）登记分期计划

{\small
\begin{longtable}{@{}L{2.4cm} L{12.9cm}@{}}
\caption{登记分期计划用例规约}\label{tab:uc-24}\\
\toprule
\textbf{标题} & \textbf{内容} \\
\midrule
\endfirsthead
\toprule
\textbf{标题} & \textbf{内容} \\
\midrule
\endhead
用例编号：用例名 & UC-24：登记分期计划 \\
\midrule
执行者 & 销售（主执行者） \\
\midrule
前置条件 & 存在成交记录，且该成交未关联分期计划 \\
\midrule
后置条件 & 新增一份状态为“未还”的分期计划，计划总额等于成交价 \\
\midrule
涉众利益 & 销售：大额成交当场可谈分期，不丢单；\newline
客户：分期减轻一次付清压力；\newline
老板：分期账目清楚，逾期可追踪 \\
\midrule
基本路径 & 1. 销售选择成交记录并填写分期期数；\newline
2. 系统校验该成交未关联分期计划；\newline
3. 系统按成交价创建分期计划（每期应还、下一期到期日、状态未还）；\newline
4. 系统反馈登记结果（含计划编号） \\
\midrule
扩展 & 2a. 已存在分期计划：2a1 系统反馈一份成交只有一份计划；2a2 用例结束 \\
\midrule
字段列表 & 分期计划＝计划编号＋成交记录＋分期期数＋每期应还＋已还金额＋下一期到期日＋状态（未还、逾期、已结清） \\
\midrule
业务规则 & ①一份成交至多一份分期计划，计划总额等于成交价；\newline ②分期期间货款按已收金额累计，不影响货品所有权转移 \\
\midrule
质量需求 & 可用性：登记一份计划不超过 3 次操作 \\
\midrule
设计约束 & 分期期数与门槛由店内规则决定，系统不设上限 \\
\midrule
\bottomrule
\end{longtable}}

\noindent（25）记录分期还款

{\small
\begin{longtable}{@{}L{2.4cm} L{12.9cm}@{}}
\caption{记录分期还款用例规约}\label{tab:uc-25}\\
\toprule
\textbf{标题} & \textbf{内容} \\
\midrule
\endfirsthead
\toprule
\textbf{标题} & \textbf{内容} \\
\midrule
\endhead
用例编号：用例名 & UC-25：记录分期还款 \\
\midrule
执行者 & 销售（主执行者） \\
\midrule
前置条件 & 存在未结清的分期计划 \\
\midrule
后置条件 & 更新分期计划的已还金额与下一期到期日，结清后状态置为已结清 \\
\midrule
涉众利益 & 客户：还款进度随时可查；\newline
销售：收款当场登记，不用事后补账；\newline
老板：逾期计划一目了然 \\
\midrule
基本路径 & 1. 销售选择分期计划并录入还款金额；\newline
2. 系统校验金额不超过未还余额；\newline
3. 系统累加已还金额并推进下一期到期日；\newline
4. 系统校验是否结清，结清则置为已结清；\newline
5. 系统反馈还款结果（含剩余应还） \\
\midrule
扩展 & 2a. 金额超过未还余额：2a1 系统反馈应按余额录入；2a2 返回1 \\
\midrule
字段列表 & 同分期计划 \\
\midrule
业务规则 & ①到期日过仍有未还即置为逾期；\newline ②还款可跨期一次结清；\newline ③结清后计划归档，不再接受还款 \\
\midrule
质量需求 & 可用性：记录一次还款不超过 2 次操作 \\
\midrule
设计约束 & 逾期状态在查询时按到期日与未还余额实时判定（超过到期日未还即逾期） \\
\midrule
\bottomrule
\end{longtable}}

\noindent（26）查看库存流水

{\small
\begin{longtable}{@{}L{2.4cm} L{12.9cm}@{}}
\caption{查看库存流水用例规约}\label{tab:uc-26}\\
\toprule
\textbf{标题} & \textbf{内容} \\
\midrule
\endfirsthead
\toprule
\textbf{标题} & \textbf{内容} \\
\midrule
\endhead
用例编号：用例名 & UC-26：查看库存流水 \\
\midrule
执行者 & 老板（主执行者）；销售、合作商户（只读） \\
\midrule
前置条件 & 无（登录即可） \\
\midrule
后置条件 & 无数据变更，返回满足条件的库存流水清单 \\
\midrule
涉众利益 & 老板：任一货品的进出有账可查，调货库存一本账；\newline
销售：交接班时核对货品状态；\newline
合作商户：对账时可当场出示流水 \\
\midrule
基本路径 & 1. 老板选择货品或时间段；\newline
2. 系统按条件查询库存流水；\newline
3. 系统返回流水清单（变更前状态、变更后状态、关联单据、经手人、时间） \\
\midrule
扩展 & 2a. 无满足条件的流水：2a1 系统反馈该条件下无记录 \\
\midrule
字段列表 & 库存流水＝货品＋变更前状态＋变更后状态＋关联单据＋经手人＋流水时间 \\
\midrule
业务规则 & ①货品每次状态变更必须写一条流水；\newline ②流水写入后不可修改、不可删除 \\
\midrule
质量需求 & 性能：单货品流水查询 3 秒内返回 \\
\midrule
设计约束 & 流水状态取值与货品状态机（07a）一致 \\
\midrule
\bottomrule
\end{longtable}}

\noindent（27）查询预约

{\small
\begin{longtable}{@{}L{2.4cm} L{12.9cm}@{}}
\caption{查询预约用例规约}\label{tab:uc-27}\\
\toprule
\textbf{标题} & \textbf{内容} \\
\midrule
\endfirsthead
\toprule
\textbf{标题} & \textbf{内容} \\
\midrule
\endhead
用例编号：用例名 & UC-27：查询预约 \\
\midrule
执行者 & 客户（无需登录） \\
\midrule
前置条件 & 存在以该手机号提交、状态为待确认且时段未过期的预约 \\
\midrule
后置条件 & 无数据变更，返回匹配的预约信息 \\
\midrule
涉众利益 & 客户：不必登录，凭手机号后四位就能找回预约编号与到店时段；\newline
销售：接待时客户能报出预约编号，核对更快；\newline
老板：查询不暴露客户完整信息 \\
\midrule
基本路径 & 1. 客户输入手机号后四位；\newline
2. 系统按后四位检索待确认且时段未过期的预约；\newline
3. 系统返回预约清单（预约编号、期望时段、货品名） \\
\midrule
扩展 & 2a. 无匹配的预约：2a1 系统反馈无预约记录 \\
\midrule
字段列表 & 查询结果＝预约编号＋期望时段＋货品名（不含姓名与手机号） \\
\midrule
业务规则 & ①按手机号后四位检索，不按姓名模糊匹配；\newline ②只返回状态为待确认且时段未过期的记录；\newline ③结果中不出现姓名、完整手机号与成交价 \\
\midrule
质量需求 & 性能：查询 2 秒内返回；可用性：输入一次提交即可查出 \\
\midrule
设计约束 & 查询不要求登录，手机号后四位只作检索条件，不作身份凭据 \\
\midrule
\bottomrule
\end{longtable}}

\noindent（28）跟进预定到店

{\small
\begin{longtable}{@{}L{2.4cm} L{12.9cm}@{}}
\caption{跟进预定到店用例规约}\label{tab:uc-28}\\
\toprule
\textbf{标题} & \textbf{内容} \\
\midrule
\endfirsthead
\toprule
\textbf{标题} & \textbf{内容} \\
\midrule
\endhead
用例编号：用例名 & UC-28：跟进预定到店 \\
\midrule
执行者 & 销售（主执行者） \\
\midrule
前置条件 & 存在约定时段在当日、状态为待到店或已联系的预定记录 \\
\midrule
后置条件 & 预定记录置为已联系、已达成或已爽约，并记下联系时间 \\
\midrule
涉众利益 & 销售：当日要联系谁、谁已到店一屏可见，不漏接待；\newline
客户：临近到店时段接到电话确认，不必白跑；\newline
老板：预定到店率可查，额度回收有据 \\
\midrule
基本路径 & 1. 销售打开当日到店名单；\newline
2. 系统列出状态为待到店或已联系的预定（姓名、手机号、货品、时段、状态）；\newline
3. 销售电话联系客户后标记已联系；\newline
4. 系统将该预定置为已联系，记下联系时间；\newline
5. 销售登记到店结论；\newline
6. 系统按结论置为已达成或已爽约，释放或保留额度；\newline
7. 系统反馈跟进结果 \\
\midrule
扩展 & 2a. 当日无待跟进的预定：2a1 系统反馈今日无待联系客户；\newline 3a. 该预定已标记过已联系：3a1 系统反馈重复标记，不再记录 \\
\midrule
字段列表 & 预定记录增记：联系时间＋到店结论（到店、未到） \\
\midrule
业务规则 & ①名单按约定时段升序排列；\newline ②标记已联系只作登记，不改变预定额度；\newline ③登记到店结论时，到店置为已达成并释放 1 条额度，未到置为已爽约，额度不释放；\newline ④已达成与已爽约的预定不再进入当日名单 \\
\midrule
质量需求 & 性能：打开名单 2 秒内返回；可用性：标记已联系与登记结论各不超过 2 次操作 \\
\midrule
设计约束 & 跟进在柜台电脑端完成；电话由销售拨打，系统不代发通知 \\
\midrule
\bottomrule
\end{longtable}}

% ==================== C 分析（承接上文，无过渡） ====================


% ============================================================
% C 阶段：分析
% ============================================================
